% =====================================================================
% preamble.tex — Configuração ABNT estrita para PPGD/Unifor
% Doutorado em Direito Constitucional — Kalyl Lamarck
% Compatível com xelatex (suporte UTF-8 + fontes nativas)
% =====================================================================

% --- Encoding e idioma ---------------------------------------------------
% xelatex le UTF-8 nativo (inputenc nao necessario)
% babel ja carregado pela classe abntex2 (option 'brazil')

% --- Fonte Times New Roman (msttcorefonts) -------------------------------
\usepackage{fontspec}
\setmainfont{Times New Roman}       % corpo
\setsansfont{Times New Roman}       % forca serifada tambem em sans
\setmonofont{Liberation Mono}        % monoespacada (em codigo)

% --- Tamanho 12pt em tudo (incluindo headings) ---------------------------
% abntex2 usa, por default, \large/\Large nos headings. Forçamos \normalsize
% para conformidade com NBR 14724 (corpo e títulos em 12pt).
\renewcommand{\ABNTEXchapterfont}{\bfseries}
\renewcommand{\ABNTEXchapterfontsize}{\normalsize}
\renewcommand{\ABNTEXsectionfont}{\bfseries}
\renewcommand{\ABNTEXsectionfontsize}{\normalsize}
\renewcommand{\ABNTEXsubsectionfont}{\bfseries}
\renewcommand{\ABNTEXsubsectionfontsize}{\normalsize}
\renewcommand{\ABNTEXsubsubsectionfont}{\bfseries}
\renewcommand{\ABNTEXsubsubsectionfontsize}{\normalsize}

% Forçar CAIXA ALTA em chapter (seção primária) — texto digitado ficará upper
% via comando \MakeUppercase no título; abntex2 já força isso por default.

% --- Espaçamento entre linhas: 1,5 (corpo); simples (notas/refs) ---------
\OnehalfSpacing

% --- Margens ABNT 3/2/3/2 cm ---------------------------------------------
\usepackage[
    a4paper,
    top=3cm,
    bottom=2cm,
    left=3cm,
    right=2cm,
    bindingoffset=0pt
]{geometry}

% --- Recuo de primeira linha: 1,5 cm -------------------------------------
\setlength{\parindent}{1.5cm}
\setlength{\parskip}{0pt}

% --- Citações: estilo autor-data ABNT (abntex2cite) ----------------------
% [alf] = autor-data ('alfabético' no jargão ABNT)
% [abnt-emphasize=bf] = título da obra em negrito
% [abnt-etal-list=0] = sem 'et al.' (lista todos os autores)
% [abnt-and-type=&] = separador entre autores no texto
\usepackage[
    alf,
    abnt-emphasize=bf,
    abnt-etal-list=0,
    abnt-and-type=&,
    abnt-thesis-year=both,
    abnt-repeated-author-omit=yes
]{abntex2cite}

% --- Hyperref sem cores (impressão acadêmica preto-e-branco) -------------
\usepackage{hyperref}
\hypersetup{
    colorlinks=false,
    pdfborder={0 0 0},
    linkcolor=black,
    citecolor=black,
    urlcolor=black,
    pdftitle={\@title},
    pdfauthor={\@author},
    pdflang=pt-BR
}

% --- Tabelas e figuras ---------------------------------------------------
\usepackage{graphicx}
\usepackage{booktabs}
\usepackage{tabularx}
\usepackage{caption}

% Legendas em TNR 12pt, alinhadas à esquerda
\captionsetup{
    font={normalsize},
    labelfont={bf},
    justification=raggedright,
    singlelinecheck=false
}

% --- Notas de rodapé: TNR 10pt, espaçamento simples ----------------------
\renewcommand{\footnotesize}{\fontsize{10pt}{12pt}\selectfont}
\setlength{\footnotesep}{14pt}        % espaço entre notas
\renewcommand{\footnoterule}{\rule{4cm}{0.4pt}}

% --- Outras configurações estéticas --------------------------------------
\setlength{\afterchapskip}{1.5\baselineskip}
\setlength{\beforechapskip}{0pt}

% --- Renomear titulo da bibliografia para REFERENCIAS (NBR 6023) ---------
% Em modo 'article' nao ha \chapter, entao redefine \bibsection
% para usar \section* em vez do \chapter* default do abntex2.
\renewcommand{\bibname}{REFERÊNCIAS}
\renewcommand{\bibsection}{%
    \section*{\bibname}%
    \addcontentsline{toc}{section}{\bibname}%
}

% --- Variáveis institucionais (usar nos textos) --------------------------
\newcommand{\universidade}{Universidade de Fortaleza}
\newcommand{\programa}{Programa de Pós-Graduação em Direito Constitucional}
\newcommand{\nivel}{Doutorado em Direito Constitucional}
\newcommand{\autorprincipal}{Kalyl Lamarck Filgueiras Estanislau Costa}
